%============================================================ 3. 문제 정의
\section{문제 정의와 평가 지표}
\label{sec:problem}

\subsection{교전 시나리오}

교전 해역 중앙에 고가치 모선 한 척이 있고, 가장자리에서 적 USV $M$척이 여러 방위로부터
동시에 접근한다. 방어정 $P$척은 적선보다 느리며, 각각 그물벽을 전개할 수 있다. 구체적인
값은 Table~\ref{tab:scenario}에 정리하였다.

% 자동 생성 — tools/make_paper_tables.py. 손으로 고치지 말 것.
\begin{table}[t]
\centering
\caption{교전 시나리오와 관측 설정. 값은 학습 체크포인트의 설정에서 직접 읽었다.}
\label{tab:scenario}
\footnotesize
\resizebox{\ifdim\width>\linewidth\linewidth\else\width\fi}{!}{%
\begin{tabular}{lrl}
\toprule
항목 & 값 & 단위 \\
\midrule
적 USV 수 $M$ & 10 & 척 \\
방어정 수 $P$ & 3 & 척 \\
적 속력 & 9.0 & m/s \\
방어정 속력 & 6.0 & m/s \\
교전 해역 한 변 & 12600 & m \\
그물벽 최대 길이 & 450 & m \\
관측 격자 반폭 & 6300 & m \\
관측 격자 해상도 & 50 & 칸 \\
적 무리 최대 수 $C$ & 4 & 개 \\
관측 픽셀 한 변 & 252 & m \\
방어정당 그물 (학습 / 평가) & 1 / 3 & 장 \\
정책 파라미터 수 & 244,389 & 개 \\
\bottomrule
\end{tabular}
}
\\[3pt]
\begin{minipage}{\linewidth}\footnotesize
방어정 속력이 적 속력보다 낮다 --- 추격이 성립하지 않으므로 접근 회랑의 선제 차단이 유일한 대응이다. 그물 보유량은 학습 시 1장, 평가·배포 시 3장이다: 정책은 결정마다 그물벽 하나를 놓으므로 결정당 행동공간이 보유량과 무관하고, 1장으로 학습한 정책이 재학습 없이 3장 설정에서 동작한다.
\end{minipage}
\end{table}


%-------------------------------------------------- 기호표
%  ★ 이 표에 기호를 추가하려면 본문에 실제로 그 기호가 있는지 먼저 확인할 것.
%    표에만 있고 본문에 없는 기호는 독자를 찾아 헤매게 만든다.
\begin{table*}[t]
\caption{주요 기호. 본 논문은 같은 문자를 두 뜻으로 쓰지 않는다 --- 특히 무리 수 $C$,
         경유점 수 $K$, 학습의 그룹 후보 수 $G$ 는 서로 다른 문자를 쓴다.}
\label{tab:nomenclature}
% 다른 표와 같은 \footnotesize 로 둔다. \small(10.95pt)로 두면 자연 폭이 커져
% \resizebox 가 발동하고, 그러면 글자가 오히려 더 작아진다(측정: 10.95 -> 9.53).
\centering\footnotesize
% 판면보다 넓을 때만 줄인다(2단 양단걸침 ≈484pt vs elsarticle preprint 390pt).
\resizebox{\ifdim\width>\linewidth\linewidth\else\width\fi}{!}{%
\begin{tabular}{@{}ll@{\hspace{1.2em}}ll@{}}
\toprule
\textbf{기호} & \textbf{뜻 (값)} & \textbf{기호} & \textbf{뜻 (값)} \\
\midrule
$M$            & 적 USV 수 (10)                       & $\mathbf{o}_p$ & 방어정 $p$ 의 자체 상태 ($\mathbb{R}^8$) \\
$P$            & 방어정 수 (3)                        & $\mathbf{q}_p$ & 방어정 $p$ 의 질의 벡터 ($\mathbb{R}^{32}$) \\
$C$            & 무리 수, 최대 (4)                    & $\mathbf{f}_i$ & 픽셀 $i$ 의 키 벡터 ($\mathbb{R}^{32}$) \\
$K$            & 방어정 1척이 지목하는 경유점 수 (2)  & $s_{p,i}$      & 방어정 $p$ 가 픽셀 $i$ 에 준 점수 \\
$G$            & 학습의 그룹 후보 수 (8)              & $(\boldsymbol{\gamma}_{\mathrm{F}},\boldsymbol{\beta}_{\mathrm{F}})$ & FiLM 배율·이동 계수 \\
\addlinespace[2pt]
$\tau_j$       & 무리 $j$ 의 위협도                   & $R_k$          & 후보 $k$ 의 팀 보상 \\
$d_{pj}$       & 방어정 $p$ $\to$ 무리 $j$ 요격점 거리 & $A_k$         & 후보 $k$ 의 그룹 상대 이득 \\
$\beta$        & 배정 유지 보너스 (식~\ref{eq:sticky}) & $\Phi$        & 위협 근접도 퍼텐셜 \\
$\Delta t$     & 층 분리 증분 (0.18)                  & $\gamma$       & 할인 계수 (0.99) \\
\addlinespace[2pt]
$\mathbf{I}$   & 요격점                               & $\mathcal{V}_k$ & $k$ 번째 선택 시점의 후보 픽셀 집합 \\
$\mathbf{t}_k$ & $k$ 번째 목표점 (연속 좌표)          & $r_{\mathrm{dup}}$ & 중복 배제 반경 (200\,m) \\
$\hat{\mathbf{r}}$ & 모선 $\to$ 요격점 단위 벡터      & $L_{\mathrm{net}}$ & 그물벽 최대 길이 (450\,m) \\
$\mathbf{a}_p$ & 방어정 $p$ 의 위치                   & $\dz$          & 대응표본 효과크기 \\
\bottomrule
\end{tabular}
}
\\[3pt]
\begin{minipage}{\linewidth}\footnotesize
격자는 $50\times50$ 이고 1\,px $=252\unit{m}$ 이다(관측 반폭 6300\,m). 결정 주기는
25\,step $=25$초이며 저수준 제어는 1초마다 돈다. 설정값 전체는
Table~\ref{tab:scenario}에 있다.
\end{minipage}
\end{table*}



방어정 속력이 적 속력보다 낮다는 점이 문제의 성격을 결정한다. 추격이 성립하지 않으므로
방어정은 적을 따라가는 대신 적이 지나갈 자리를 예측하여 미리 점유해야 한다. 그물벽은 두
좌표를 잇는 선분으로 정의되며, 방어정이 그 선분을 따라 이동하면서 전개한다. 전개된 그물은
해면에 고정되지 않고 표류하므로, 적에 가깝게 전개할수록 유효 시간이 길다.


적 편성은 세 가지 공격 포메이션으로 구성한다.
\begin{itemize}[leftmargin=*, itemsep=3pt, topsep=3pt]
\item \textbf{집중(concentrated)}: 적 10척이 한 방위에 모여 동시에 접근한다. 방어정 3척이
      모두 같은 방향을 막으면 되므로 배정 판단의 여지가 작다.
\item \textbf{양동(diversionary)}: 적이 세 방위에서 $4/3/3$ 으로 나뉘어 \emph{동시에}
      도래한다(기준 방위로부터 약 $130^\circ$, $230^\circ$ 이격). 방어정 수와 무리 수가
      같아 \emph{여유가 전혀 없으므로}, 어느 방어정을 어느 무리에 붙이는지와 그 경로가
      교차하지 않는지가 결정적이다.
\item \textbf{파상(wave)}: 한 방위에서 3개 단(rank)이 거리차를 두고 줄지어 도달한다
      (최근접 4{,}000\unit{m}, 단 간격 1{,}000\unit{m}). 앞 단을 막은 자원이 뒤 단을
      막을 수 있어야 하므로 \emph{어느 자리에} 그물벽을 놓는지가 다음 단의 차단 가능성을
      좌우한다.
\end{itemize}

세 포메이션을 모두 보고하는 이유는 어느 하나만 쓰면 결론이 포메이션에 종속되기 때문이다.
특히 집중 포메이션은 \S\ref{sec:results}에서 보이듯 천장 효과가 있어 단독으로는 어떤 방법도
변별하지 못한다.

\subsection{평가 지표}
\label{sec:metrics}

방어 성과의 1차 지표는 포획률이다.
\begin{equation}
\capr = \frac{N_\mathrm{cap}}{N_\mathrm{cap} + N_\mathrm{brc}}
\label{eq:capture_rate}
\end{equation}
여기서 $N_\mathrm{cap}$과 $N_\mathrm{brc}$는 각각 포획된 적선 수와 모선까지 돌파한 적선
수이다.

그러나 포획률 하나로는 방어의 질을 판정할 수 없다. ``그물을 퍼부어 올린 포획률''과 ``적은
자원으로 얻은 포획률''이 구분되지 않고, 방어정끼리 충돌하여 잃은 손실도 드러나지 않는다.
따라서 다섯 관점에서 총 10개 지표를 정의하고 모두 동일한 통계 절차로 보고한다
(Table~\ref{tab:metrics}).

\begin{table*}[t]
\centering
\caption{평가 지표 10종. 방향은 개선이 어느 쪽인지를 나타낸다.}
\label{tab:metrics}
\footnotesize
\begin{tabular}{llcp{0.44\linewidth}}
\toprule
관점 & 지표 & 방향 & 정의 \\
\midrule
\multirow{2}{*}{방어 성과}
 & 포획률           & $\uparrow$ & 식~\eqref{eq:capture_rate} \\
 & 돌파 수          & $\downarrow$ & 모선까지 도달한 적선 수 \\
\midrule
\multirow{2}{*}{안전}
 & 충돌률           & $\downarrow$ & 아군 상호 및 모선 충돌 건수(척당) \\
 & 그물 걸림        & $\downarrow$ & 아군이 설치된 그물에 얽힌 횟수 \\
\midrule
\multirow{2}{*}{자원 효율}
 & 포획당 그물      & $\downarrow$ & 적 1척 포획에 소모한 그물 수 \\
 & 포획당 이동거리  & $\downarrow$ & 적 1척 포획당 팀 총 이동거리 \\
\midrule
기동 비용
 & 총 선회량        & $\downarrow$ & 누적 방위 변화량(궤적 부드러움) \\
\midrule
\multirow{3}{*}{교전 품질}
 & 포획 이격거리    & $\uparrow$ & 포획 지점의 모선 이격 --- 공간축 조기 제압 \\
 & 평균 포획 시각   & $\downarrow$ & 포획이 일어난 시각 --- 시간축 조기 제압 \\
 & 전개 시 적 거리  & $\downarrow$ & 그물 전개를 시작한 순간 최근접 적까지 거리 \\
\bottomrule
\end{tabular}
\\[3pt]
\begin{minipage}{\linewidth}\footnotesize
포획당 그물·포획당 이동거리처럼 성과로 나눈 지표를 쓰는 이유는, 총량만 보면 ``아무것도
하지 않은 조건''이 가장 효율적으로 보이기 때문이다. 조기 제압을 거리축과 시간축 양쪽으로
재는 이유는 멀리서 잡아도 늦게 잡으면 다음 파를 막지 못하고, 빨리 잡아도 모선 코앞이면
여유가 없기 때문이다.
\end{minipage}
\end{table*}

전개 시 적 거리를 낮을수록 좋은 지표로 두는 것은 그물의 물리적 성질 때문이다. 그물은 해면에
고정되지 않고 표류·확산하므로, 적에서 멀리 떨어진 곳에 미리 뿌린 그물은 적이 도달할 무렵
이미 형태를 잃는다.
